\subsection{变换（克莱因–李，1869-1872）。}


1860 年前后，有限集的置换群理论正在发展，并开始得到应用（塞雷、克罗内克、马蒂厄、若尔当）。另一方面，当时正处于全盛时期的不变量理论，使数学家们熟悉了某些对合成封闭的几何变换的无穷集合（特别是线性变换或射影变换）。但是，在若尔当 1868 年关于“运动群”（三维欧几里得空间中位移群的闭子群）的工作{[}174 b{]}之前，这两股思想潮流之间似乎还没有建立起任何自觉的联系。

1869 年，年轻的费利克斯·克莱因（1849-1925），普吕克的学生，在柏林与年长他几岁的挪威人索弗斯·李（1842-1899）结为好友，把他们聚到一起的是两人对普吕克的“直线几何学”、特别是直线丛理论的共同兴趣（第~133 页）。正是在这个时期前后，李形成了他最富原创性的想法之一：把不变量的概念引进分析与微分几何；其来源之一是他观察到，微分方程“用求积法”积分的经典方法全都依赖于这样的事实：方程对于一个“连续”的变换族是不变的。李使用这一想法的第一批工作（由克莱因整理）正是从 1869 年开始的；他在其中研究了“雷耶丛”（在与四面体的各面相交于具有给定交比的 4 个点的直线所成的集）以及以该丛的直线为切线的曲线与曲面（{[}202{]}，vol.~I, Abh. V, pp.~68-77）：他的方法依赖于雷耶丛在保持四面体各顶点不变的三参数交换群（PGL(4, C) 的一个极大环面）下的不变性。同一想法也支配着克莱因与李于 1870 年春同在巴黎时合写的著作（{[}182{]}，v. I, pp.~416-420）；他们在其中实质上确定了平面的射影群 PGL(3, C) 的连通交换子群，并研究了其轨道的几何性质（以 V 曲线或 V 曲面之名）；这使他们通过一个统一的程序，得到各种各样的曲线——代数的或超越的——的性质，例如 \(y = cx^m\) 或对数螺线。两人的证词一致强调伽罗瓦与若尔当的理论给他们留下的深刻印象（若尔当关于伽罗瓦的评注已于 1869 年发表在 Math. Annalen 上；无论如何，李早在 1863 年就已听说过伽罗瓦理论）。1871 年开始对非欧几何感兴趣的克莱因，在其中看到了他为一切已知几何学寻找一个分类原理的研究的开端，这项研究将在 1872 年把他引向“埃尔朗根纲领”。至于李，他在 1873 年给 A. 迈尔的一封信（{[}202{]}，vol.~5, p.~584）中，把他关于变换群的思想的起源定于他旅居巴黎期间，而在 1871 年的一部著作（{[}202{]}，vol.~I, Abh. XII, pp.~153-214）中，他已经使用了“变换群”这一术语，并明确提出了确定 GL(n, C) 的所有子群（“连续的或不连续的”）的问题。说实在的，克莱因与李两人想必都曾难以深入这个新的数学宇宙，克莱因谈到新问世的若尔当的《专题论文集》时，说它是一本“上了七道封印的书”（{[}182{]}，v. I, p.~51）；他在别处谈到（{[}182{]}，v. I, pp.~424-459）时写道：“凡涉及连续算子群的启发式观念的一切，特别是凡涉及微分方程或偏微分方程的积分的一切，都属于李。他后来在他的连续群理论中发展起来的一切概念，当时都已在他那里以萌芽的状态存在，然而却加工得如此之少，以致有许多细节我都不得不说服他，例如甚至在开头，连 V 曲线的存在也是在这样的（说服）之下，即在漫长的交谈过程中（才被接受的）”（{[}182{]}，v. I, p.~415）。

